\acknowledgements

To be honest, I would rather finish the acknowledgements after the defense is done because I view it as a means to celebrate, a silent speech to address the long journey of Ph.D. I have partaken. However, the standard practice says otherwise, so I have to write it down now, before sending the thesis to the committee. I reckon this is also a good thing, as it adds more humbleness to the overall tone, something we all need to be reminded of from time to time.

I would first like to express my deepest gratitude to my advisor Antonio (Tuca) Auffinger for his guidance and support throughout these years, as well as the other two committee members Elton Hsu and Sandy Zabell. Such gratitude is much deeper than the ordinary, or dare I say, routine acknowledgements that would usually appear in a Ph.D. thesis. Allow me to elaborate.

During my third year in the graduate program at Northwestern, I suffered from a serious sickness that prevented me from conducting research work for almost two years. What made things worse was that period of time overlapped with the peak of the COVID-19 pandemic. I had to pick up new research problems from scratch after getting out of the hospital at the end of 2021. The rehabilitation process already took most of my energy, and my empty research portfolio told me that my path to academia was no more. I needed time, a lot of time, to put my health, my passion, and my career back together from all the broken pieces scattered across my life. Time was never on my side.

At this time, instead of telling me my time was up with the Ph.D. program, Tuca gave me more time and space to work on my problems, find industry jobs, and most importantly, to think, reflect, and heal. He told me that intelligence and talent were not the problem, but my state of mind was. He supported the worst version of me to become, in my opinion, the current best version of me. Elton also provided me with a lot of help and mental comfort. The math department at Northwestern chose to provide me not only with all the excellent resources in mathematics but also a sanctuary for lost young folks. This place will forever be a beacon of hope, a place that reminds me of the kindness and intelligence in human beings, and a source of solace during the darkest times.

With the above being said, please bear with me for some extra words for other important people in my life who backed me up along the way.

First, I would like to thank my best friend Xiaoyu Li for his life-long support and for taking care of me during my surgery periods. He generously shared his place with me and looked after me when I needed to constantly go in and out of the hospital due to a lack of patient rooms during the pandemic. Interestingly, he was also the person who drove across the country from Houston to Evanston to move me and my belongings at the beginning of my Ph.D. program. It is my utmost honor and privilege to have such a friend whom I can trust with my life and soul.

Then, I would like to thank my collaborators and dear friends Zichao Wang and Nan Jiang for their inspiration and guidance on the research papers we published. Not only did they offer me their wisdom in computer science, but also their invaluable friendship. Without them, the completion of this thesis would not have been possible.

Next, I would like to express my gratitude to my colleague Chenyao Zhang at Google, who mentored me during my internship. She was a good mentor and friend who not only guided me through my time at Google but also shared her rich industry experience with me. She was the one who pushed for my hiring decision and the reason I had such an amazing team at Google amidst mass layoffs and hiring freezes. Though not knowing me for long, she was willing to trust and support a stranger without hesitation. I will strive to do the same and pass down the kindness to others, just as she did.

Moreover, I would like to thank my good friends Xiaoyini Zhang and Zeyu Li for offering me their places to stay during my internship in the Bay Area. When I received my offer at Google, I was still busy with my rehabilitation process, and relocation across the States was a huge challenge for me. They immediately decided to share their places with me despite their own stressful lives. They, like all the people I have mentioned, are the ones who made me believe in the goodness of mankind—a reminder for me to be kind regardless of the circumstances.

In addition, I would like to thank my violin teacher Addison Teng for his care and inspiration during my hard time of rehabilitation. I met Addison three years ago when I felt that I needed some extra help to not give up the violin, the instrument which I have played for more than 25 years, since I was 4. Addison was one of those teachers who not only excel at what they teach but also care deeply about the well-being of their students. He taught me to reflect on my own internal shortcomings and to face hardship without fear. He taught me to be a better musician and a better human. Playing the violin is an everlasting service to oneself, a procedure that teaches us that we should always strive for the best and never compromise on our core values and standards. I am fortunate to have him by my side.

There are many, many others to whom I would like to express my appreciation, such as my alumni and friends Yuxin Zhou and Mengxuan Yang for their fruitful discussions regarding my research problems. Others include my old friends Anqi Sun and her partner Ke Jin, who invited me on their travels during my stressful times. The list goes on and on.

Finally, I would like to thank my family, my parents for their unconditional love and support. People say that one's family shapes one's personality. There is no doubt that my life was shaped by love, warmth, and wisdom. My parents passed down the most valuable things in our civilization that will shine through the course of history like a burning star. Now it is my turn to uphold their legacy and make my contribution to this world as a mathematician, a scientist, and most importantly, a person with curiosity, hope, and love.